\chapter{Of Labor, Shadowing, and the Assignment}
\label{art:labor}

\articlenote{In which is set down how the silo's citizens are matched to their work, how they are taught to do it, and the principle by which some labor is chosen freely and some is assigned by need.}

\section{The Charge of Labor and the Principle of Assignment}

Every citizen, from the sixteenth naming-day onward, is bound to contribute labor to the silo's keeping, in such a trade or profession as the citizen's own aptitude, the silo's needs, and the Pact's requirements together determine. The silo contains one hundred and forty-four levels, built and provisioned to hold many thousands of souls; the work required to keep so many fed, watered, warmed, aired, and sheltered is immense and various. No citizen's labor is purchased by chits or offered freely; it is owed to the silo itself as a debt of citizenship.

\begin{clauses}
\item A citizen may petition to enter a trade or profession for which the citizen has aptitude and desire, and such petition shall not be refused save where the silo's need for that trade is already met and other work goes short. The silo's welfare takes precedence over the individual citizen's preference, but a citizen's own petition is weighed first.
\item Where a trade is in short supply and the silo's need is critical, the Mayor and the Head of the relevant Department may jointly assign a citizen to that trade, without the citizen's consent, provided Judicial confirms that the assignment is not made in punishment under Article 17.
\item A citizen assigned to a trade without consent shall be tested for fitness in that trade; a citizen found unfit may petition Judicial for reassignment, and Judicial's finding shall bind.
\item A citizen who neglects the labor assigned, or who diverts the silo's tools, materials, or the citizen's own owed labor-time to a private study, pursuit, or fixation of the citizen's own devising, wastes what the silo cannot spare. The Head of the citizen's Department shall first correct such neglect by closer oversight or reassignment to labor better suited to occupy the citizen fully; where the neglect persists notwithstanding, the Head shall refer the matter to Judicial.
\end{clauses}

\section{The Testing of Aptitude}

From a citizen's twelfth naming-day forward, the schools under Article 15 shall assess each child's aptitude for various kinds of labor: the working of tools, the tending of living things, the care of people, the working with numbers, the keeping of records, and such other aptitudes as the silo's work requires. A citizen reaching majority shall have some knowledge of which trades might suit the citizen's own nature.

\begin{clauses}
\item The testing of aptitude is not a trial to be passed or failed; it is a guidance, carried into the citizen's majority and weighed when the citizen first petitions for entry into a trade.
\item A citizen may petition to enter a trade notwithstanding the aptitude testing's results, but a Head of a Department may weigh the results against the petition: where a citizen's expressed desire runs sharply against demonstrated aptitude, the Head may require an extended trial period or additional testing before accepting the citizen's shadow.
\item A citizen who petitions for a trade for which the silo has no current need, but for which the citizen has clear aptitude, may be placed on a waiting list; if the silo's need for that trade grows, the waiting citizen has first claim when a shadow becomes available.
\end{clauses}

\section{The Three Tiers of Shadowing and Labor Entry}

The silo's work is divided into three categories by the length and rigor of training required. A citizen enters a trade through one of three pathways: a full shadowing of years spent learning from a departing master; a short shadowing of weeks or months for moderately skilled labor; or direct assignment to a simple task requiring little or no formal training.

\begin{clauses}
\item Full shadowing is the entry into trades of high skill, significant responsibility, or access to restricted knowledge or dangerous instruments. The shadowing shall not be less than three years nor more than seven, at the sole determination of the master being shadowed; the shadow shall live and work alongside the master, learning not only the technical work but the judgment and reliability the trade requires.
\item Short shadowing is the entry into trades of moderate skill or administrative responsibility. The shadowing shall be no fewer than four weeks and ordinarily no longer than one season, or two seasons in exceptional cases, at the determination of the master or the Head of the Department.
\item Direct assignment without formal shadowing is permitted for labor of little skill requirement, such as simple carrying, sorting, distribution, or maintenance tasks. A citizen assigned to such work receives instruction as the work requires, but no formal apprenticeship period.
\item A citizen who has shadowed under a master may not be removed from that shadow by the citizen's own request, nor by the master's, save with the consent of the other party or a finding by Judicial that the shadow is unfit. A shadow abandoned by either party without cause is answerable under Article 17 for breach of a commitment made before Judicial.
\end{clauses}

\section{Mundane Labor and the Simple Trades}

The silo requires much work of simple character: the carrying of goods and messages, the distribution of supplies, the painting and repair of common surfaces, the sorting of salvage, the cleaning of common areas, the tending of the simple tasks of kitchen and cafeteria. These trades require no formal shadowing, little specialized training, and are open to any citizen of sufficient strength or steadiness for the work.

\begin{clauses}
\item A citizen wishing to enter such work may do so upon the petition to the relevant Head of Department or, for work not assigned to any Department, to the Mayor. The Head shall assign the citizen to the work or place the citizen on a list of those available.
\item No license, restricted access, or formal testing is required for such work, save that a citizen assigned to the keeping or moving of the silo's stored goods shall be spot-checked for honesty by periodic audits of what is carried and what is delivered.
\item A citizen engaged in simple labor may be reassigned by the Head of the Department at will, with one cycle of ten days' notice.
\end{clauses}

\section{Skilled Trades and Full Shadowing}

The silo's vital functions depend on trades of genuine skill: the tending of the generator and the works of water and air, the working of metal and the repair of machinery, the keeping of food stocks and the growing of crops, the healing of the sick, the keeping of records, the investigation of crime and the carrying out of sentences, the teaching of the young. These trades require years of learning under a master and may not be entered save through full shadowing.

\begin{clauses}
\item A citizen wishing to enter a skilled trade shall petition the Head of the relevant Department. The Head shall examine the citizen's aptitude testing and interview the citizen, and may accept or refuse the petition. A refusal may be appealed to Judicial.
\item Where a shadow is accepted, a master who possesses the trade and plans to step down from it shall take the shadow as an apprentice. A master who does not plan to step down may be required by the Head of Department to take a shadow if the Department's needs require it, provided Judicial finds the requirement just.
\item The shadow shall be tested, formally and informally, throughout the shadowing period. A shadow found unfit may be dismissed by the master or the Head of Department, provided Judicial is informed; a shadow dismissed may petition Judicial for reassignment to a different trade.
\item Upon the completion of shadowing, the master shall attest the citizen's readiness to the Head of Department. The Head shall confirm the citizen in the trade, or, if the Head finds the citizen unfit despite the master's attestation, shall refer the matter to Judicial for a ruling.
\end{clauses}

\section{Critical Roles and Order-Adjacent Trades}

Certain trades carry access to knowledge or authority that goes beyond the ordinary silo's work: the technicians of Information Technology, the Sensor and Wallscreen workers, the Cleaning-Lab technicians, the teachers of the young, the Officers of Judicial, the Physicians of the Infirmary. These roles require not only full shadowing but additional testing and vetting, beyond what the shadowing system alone determines.

\begin{clauses}
\item A citizen entering a critical role shall undergo, in addition to the usual shadowing, such further testing and vetting as the relevant Head of Department determines. This vetting may include examination by Judicial, investigation by the Officers of Judicial, and assessment of the citizen's steadiness, discretion, and loyalty to the Pact.
\item No citizen shall be compelled to undergo such vetting, and a citizen refusing it shall not be forced into a critical role; a citizen may withdraw from a shadow for a critical role at any point, and shall not be punished for the withdrawal.
\item A citizen confirmed in a critical role remains subject to this vetting throughout the citizen's tenure; a Head of Department may require additional vetting at any time, and may remove a citizen from a critical role if the vetting reveals cause.
\item The specifics of the vetting for any particular critical role are not set down in this Pact, but are determined by the Head of Department and are kept in the Judicial archive. No vetting procedure may contradict this Pact or require a citizen to violate it.
\item No citizen afflicted with the Syndrome may enter, hold, or continue in any critical role, nor may any citizen afflicted with the Syndrome be confirmed as a Head of Department. Any citizen confirmed in such a role who falls afflicted with the Syndrome shall immediately report the affliction and resign from the role, and the relevant Department or the Mayor shall proceed with the shadowing and vetting process to confirm a successor.
\end{clauses}

\section{The Management of Labor Supply and Citizen Assignment}

The silo's various Departments require different numbers of citizens in different trades, and these needs change as the silo ages, as citizens retire or are removed by death or cleaning, and as the work itself changes. Where the silo's need for a particular trade cannot be met by citizens' own petitions, the system of assignment ensures adequate staffing.

\begin{clauses}
\item Each Head of Department shall, annually or as the Department's needs change, report to the Mayor the Department's required staffing levels by trade. The Mayor, in consultation with the Heads of Department and the Assembly under Article 20, shall establish the silo's labor needs for the coming year.
\item Where citizens' own petitions for a trade exceed the silo's need for that trade, the excess petitioners are placed on waiting lists by their petition date; as citizens retire or die, the waiting list is filled in order.
\item Where the silo's need for a trade exceeds the number of citizens who have petitioned for it, and waiting lists are exhausted, the Mayor and the relevant Head of Department may jointly assign citizens to that trade, subject to Judicial confirmation as provided in Section 1.
\item No citizen of majority shall be left without assignment to some trade or labor, and no citizen's trade shall be left vacant without a shadow being prepared to fill it, save in genuine emergency. A citizen unable to find a trade to petition shall be assigned to mundane labor by the Mayor, pending opportunity for training in a more specialized field.
\end{clauses}

\section{Reassignment, Retirement, and the Management of Change}

A citizen's first assigned trade is not permanent. Citizens may petition for reassignment to a different trade after a minimum period of service (five years for skilled trades, one year for simple labor); citizens may be reassigned by their Department if the citizen's fitness is questioned or the silo's needs change; citizens may retire from labor upon reaching such age as the Infirmary and the Mayor jointly determine.

\begin{clauses}
\item A citizen petitioning for reassignment to a different trade shall undergo the usual petition and shadowing process for the new trade, except that a citizen's prior aptitude testing may be waived by the Head of the new Department if the citizen's years of successful service in another skilled trade demonstrate the necessary aptitude.
\item A citizen removed from a trade for cause (unfitness, abandonment of the shadowing, or breach of the citizen's duties) may petition for reassignment to a different trade, and Judicial shall rule on the petition; Judicial may grant reassignment, deny it, or require a period of service in mundane labor before permitting a reassignment to a new skilled trade.
\item A citizen reaching such age as the Infirmary judges the citizen no longer fit for the citizen's current trade may be reassigned to simpler labor, or, if unfit for any labor, may petition for retirement. Retirement ends the citizen's labor obligation under this Article, and the citizen becomes the responsibility of the silo under Article 19.
\item Where a Department's need for a particular trade diminishes, and the citizens holding that trade cannot be reassigned within the Department, they shall be reassigned to other Departments' needs, or retrained for related work, or retired if age or fitness permits. No citizen shall be removed from employment without alternative assignment or retirement.
\end{clauses}

\section{The Rights and Duties of the Shadowed}

A citizen in shadowing is in a status between ordinary citizenship and full membership in a trade. The shadow owes the master and the silo earnest labor and obedience to the master's teaching; the master, the Department, and the silo owe the shadow fair treatment, protection from harm, and assessment according to genuine fitness.

\begin{clauses}
\item A shadow shall receive the same ration as any other citizen under Article 11, and shall not be docked for tools, materials, or lodging during the shadowing period, save where the shadow has been assigned to lodging within the Department's own quarters, as a Down Deep trade requiring residence in Mechanical's own space.
\item A shadow may not be punished by the master for ordinary mistakes in learning, but may be dismissed for willful negligence, refusal to obey instruction, or breach of the Pact. A shadow dismissed shall be paid for the time worked and shall not be bound by any contract to repay training costs.
\item A master may not strike, deprive of food, or otherwise assault a shadow, and a shadow subjected to such treatment may appeal to Judicial or the Sheriff under Article 7 for relief and may be released from the shadowing without penalty.
\end{clauses}
